\documentclass[french]{resume}

\begin{document}

\cvheader{Ahmed Guetti}
  {Lille, France \sep +33 6 64 77 86 99 \sep
   \href{mailto:ahmedguetti12@gmail.com}{ahmedguetti12@gmail.com} \sep
   \href{https://github.com/AhmedGuetti}{github.com/AhmedGuetti} \sep
   \href{https://linkedin.com/in/ahmed-guetti-09552b1b6}{linkedin.com/in/ahmed-guetti-09552b1b6}}

% ---------- FORMATION ----------
\section{Formation}
\cventry{Master 2 Intelligence Artificielle (IAFA), parcours Informatique Graphique et Image ; classement : 16/43}
  {Université Toulouse III -- Paul Sabatier}{Sept.\ 2026}{Toulouse, France}
\begin{itemize}
  \item Travail d'initiation à la recherche (TIR) : « Apprentissage automatique pour l'optimisation combinatoire » (TSP, branch-and-bound pour MILP, SAT), co-écrit avec deux camarades ; présenté sous forme de poster
  \item Cours principaux : intelligence artificielle, planification, rendu 3D temps réel (OpenGL), géométrie et traitement de maillages, vision par ordinateur, systèmes répartis, bases de données réparties
  \item Projets de cours : ré-identification de personnes (LOMO + XQDA) ; moteur de rendu glTF PBR (Cook--Torrance, ombres PCF) ; animation squelettique et skinning (linear blend skinning CPU vs.\ GPU)
\end{itemize}

% ---------- EXPÉRIENCE ----------
\section{Expérience}
\cventry{Ingénieur IA/ML et ingénieur QA (alternance)}
  {Magellan (ex-Worldline MeTS / M3S)}{Oct.\ 2025 -- sept.\ 2026}{Seclin, France}
\begin{itemize}
  \item Systèmes agentiques de niveau production sur des projets IA internes, en parallèle d'un rôle QA/tests d'API (RATP).
  \item \textbf{BugPredict :} conception d'un pipeline multi-agents de détection de bugs avec LangGraph ; refonte d'un monolithe en architecture modulaire (système de plugins ToolRegistry, périmètre de recherche basé sur Pydantic, arêtes conditionnelles dans le graphe).
  \item Évaluation de BugPredict sur BigVul (précision 1,0, rappel 0,235, F1 0,381) ; tableau de bord et interface PySide6 suivant le pipeline en direct.
  \item Mise en place de la CI/CD GitLab (étapes prepare, build, release, notify) avec publication automatique d'artefacts versionnés.
  \item \textbf{ARGUS :} moteur d'agents générique configurable en YAML, avec orchestrateurs interchangeables (Gemini, LangGraph) et registre d'outils par plugins ; éditeur visuel de configuration (React, React Flow, Zustand).
  \item Évaluation d'un agent de détection d'incohérences spécification/code sur CASCADE, des sous-ensembles de SWE-bench et des jeux de données oracles synthétiques.
  \item Supports de démo et de passation (documentation, schémas, vidéo) pour les deux projets.
\end{itemize}

\cventry{Stagiaire de recherche}{ISIS Champollion, équipe CHART}{Avr.\ -- juil.\ 2025}{Castres, France}
\begin{itemize}
  \item \textbf{EasyTalk :} extraction d'événements et d'expressions temporelles dans des textes médicaux français (HeidelTime, Mistral, GLiNER, CamemBERT), servie via FastAPI, annotation sous INCEpTION.
  \item \textbf{EasyRMM :} refonte d'une application web existante.
\end{itemize}

\cventry{Stagiaire de recherche}{CNRS UMR LISA}{Avr.\ -- mai 2025}{Corte, France}
\begin{itemize}
  \item Co-développement d'un système de traduction automatique par TAL pour la langue corse.
  \item Pipeline de scripts AWK pour l'analyse statistique automatique des sorties du modèle, permettant d'identifier les erreurs grammaticales récurrentes.
  \item Amélioration d'environ 10\,\% de la qualité de traduction sur les phrases complexes grâce à un réglage ciblé des hyperparamètres.
\end{itemize}

% ---------- PROJETS ----------
\section{Projets}
\cventry{Bibliothèque mathématique C++ (namespace \texttt{gm}) et algorithmes de ML from scratch}
  {gmath / Machine Learning en C++}{2026 -- aujourd'hui}{Projet personnel}
\begin{itemize}
  \item Implémentation d'algorithmes de ML from scratch en C++ (1-NN/k-NN) sur une bibliothèque mathématique statique réutilisable (fonctions de distance, KD-tree générique utilisé pour le ML et le rendu).
\end{itemize}

\cventry{Plateforme d'organisation des candidatures}{Projet Ninurta}{2026 -- aujourd'hui}{Projet personnel}
\begin{itemize}
  \item Conception d'un système multi-agents qui adapte mon CV à chaque offre à partir de mon expérience.
  \item Statistiques sur les candidatures envoyées et suggestions pour ajuster leur volume et leur ciblage.
\end{itemize}

% ---------- COMPÉTENCES ----------
\section{Compétences et centres d'intérêt}
\begin{cvskills}
  \cvskill{Langages}{Python, C++, Rust, TypeScript/JavaScript, SQL}
  \cvskill{IA/ML}{LangGraph, LangChain, Pydantic, systèmes agentiques et évaluation, TAL (CamemBERT, GLiNER, API de LLM), vision par ordinateur}
  \cvskill{Outils}{OpenGL, FastAPI, React, PySide6, GCP, GitLab CI/CD, Git, WSL2/Linux, Neovim}
  \cvskill{Langues}{arabe (langue maternelle), français, anglais}
\end{cvskills}

\end{document}
